% T3.1 / T3.2 — Section IV-B (methodology), IV-C (precision), IV-D (long context).
% Closes R4-10 (every methodology item listed), R2-4 (soften the optimal-F
% claim) and R4-11 (long-context evaluation), and carries the T3.4 result.
%
% Replaces the submitted Section IV-C and Table IV wholesale. The submitted
% numbers were produced without a working angle encoder and without modelling
% datapath saturation, so none of them describe the synthesised hardware.
%
% All figures below are measured: full WikiText-2 test split, fp16, ctx 2048
% unless stated. Nothing here is a placeholder.

\subsection{Evaluation Methodology}
\label{sec:evalmethod}

Hardware behaviour is injected into the models by replacing the rotary
embedding with a bit-accurate model of the synthesised datapath --- the same
encoder, range reduction, block normalisation and shift-add pipeline described
in Section~\ref{sec:arch}, verified against the RTL to $0.002^{\circ}$. Only
the \emph{rotation} is approximated; weights, activations and the queries and
keys themselves remain in \texttt{float16} throughout, so the numbers isolate
the cost of the rotation hardware rather than of quantisation in general.
Table~\ref{tab:method} records the protocol in full.

\begin{table}[t]
\caption{Evaluation protocol.}
\label{tab:method}
\centering
\begin{tabular}{ll}
\toprule
Dataset & WikiText-2 (raw), test split, joined \\
Tokens scored & 341\,302 (LLaMA-2-7B), 298\,931 (Qwen2-7B) \\
Context length & 2048 (1024 / 2048 / 4096 in \S\ref{sec:longctx}) \\
Stride & equal to context; each token scored once \\
Batch size & 1 \\
KV cache & disabled \\
Model precision & \texttt{float16} \\
Approximated & RoPE rotation only \\
$\mathbf{q},\mathbf{k}$ precision & \texttt{float16}, unquantised \\
Loss & explicit shift, $\text{logits}_{:,:-1}$ vs.\ $\text{input}_{:,1:}$ \\
Harness & custom; \emph{not} lm-eval-harness \\
Normalisation & shared exponent per (head, token) \\
Framework & transformers 5.17.0, torch 2.10.0 \\
\bottomrule
\end{tabular}
\end{table}

Two points are worth stating plainly. First, perplexity is computed by an
explicit shift rather than by passing \texttt{labels} to the model, because
the framework applies its own shift internally and the two compose silently
into an off-by-one evaluation. Second, this is a custom harness rather than a
standard one, so the absolute baselines should be compared against the
floating-point row of the same table rather than against published figures.

\subsection{Precision Requirement}
\label{sec:precision}

Table~\ref{tab:fsweep} sweeps the fractional width $F$ end-to-end.

\begin{table}[t]
\caption{Perplexity versus fractional width, Binary UD ($N=10$),
WikiText-2, $\mathrm{ctx}=2048$. $\Delta$ is relative to the
\texttt{float16} baseline of the same model.}
\label{tab:fsweep}
\centering
\begin{tabular}{ccccc}
\toprule
& \multicolumn{2}{c}{LLaMA-2-7B} & \multicolumn{2}{c}{Qwen2-7B} \\
\cmidrule(lr){2-3}\cmidrule(lr){4-5}
$F$ & PPL & $\Delta$ & PPL & $\Delta$ \\
\midrule
baseline & 5.4680 & ---        & 7.1405  & --- \\
14       & 5.4681 & $+0.001\%$ & 7.1405  & $-0.000\%$ \\
13       & 5.4681 & $+0.001\%$ & 7.1403  & $-0.002\%$ \\
12       & 5.4682 & $+0.003\%$ & 7.1408  & $+0.005\%$ \\
11       & 5.4683 & $+0.005\%$ & 7.1417  & $+0.018\%$ \\
10       & 5.4691 & $+0.020\%$ & 7.1478  & $+0.103\%$ \\
9        & 5.4723 & $+0.078\%$ & 7.5943  & $+6.36\%$ \\
8        & 5.4854 & $+0.318\%$ & 16.527  & $+131\%$ \\
7        & 5.5498 & $+1.50\%$  & 583.71  & $+8075\%$ \\
6        & 6.0984 & $+11.5\%$  & 3407.5  & $+47621\%$ \\
\bottomrule
\end{tabular}
\end{table}

The two models do not agree, and the disagreement is the substantive result.
LLaMA-2-7B is within $0.32\%$ of its baseline at $F=8$ and degrades gracefully
below that. Qwen2-7B has already lost $6.4\%$ at $F=9$ and collapses at $F=8$
--- two orders of magnitude of perplexity across a single bit. A precision
adequate for one model is unusable for the other, on identical hardware.

The cause is the frequency schedule, not model scale: both are 7B parameter
models with the same head dimension. Qwen2 uses a RoPE base of $10^{6}$
against LLaMA-2's $10^{4}$, so its $\Theta_i = b^{-2i/d}$ spans a far wider
dynamic range and the low-frequency dimension pairs are assigned angles small
enough that coarse fractional widths cannot resolve them at all. Once those
angles fall below the datapath LSB the affected pairs stop rotating, and the
positional signal they carry is lost rather than merely degraded --- which is
what produces a cliff rather than a slope.

We therefore do not report a single optimal precision. \textbf{$F=11$ is the
defensible recommendation}: $+0.005\%$ and $+0.018\%$ on the two models, one
bit of margin above the point at which Qwen2 begins to move. $F=10$ remains
within $0.11\%$ on both and is the aggressive choice where area dominates.
Designs targeting a known model may go further --- LLaMA-2 tolerates $F=8$,
saving three bits of datapath --- but a general-purpose accelerator should
provision for the wider base, since recent models have trended towards larger
RoPE bases to support longer contexts.

\subsection{Long-Context Behaviour}
\label{sec:longctx}

Because $\theta = m\Theta_i$ grows without bound in the token position, a
natural concern is that rotation error accumulates with sequence length.
Table~\ref{tab:longctx} shows it does not.

\begin{table}[t]
\caption{$\Delta$ perplexity versus context length, Binary UD ($N=10$).}
\label{tab:longctx}
\centering
\begin{tabular}{llccc}
\toprule
Model & $F$ & 1024 & 2048 & 4096 \\
\midrule
LLaMA-2-7B & 10 & $+0.026\%$ & $+0.020\%$ & $+0.020\%$ \\
LLaMA-2-7B & 9  & $+0.098\%$ & $+0.078\%$ & $+0.086\%$ \\
Qwen2-7B   & 10 & $+0.146\%$ & $+0.103\%$ & $+0.080\%$ \\
Qwen2-7B   & 9  & $+3.05\%$  & $+6.36\%$  & $+4.39\%$ \\
\bottomrule
\end{tabular}
\end{table}

At a precision the model tolerates, the relative degradation is flat or mildly
decreasing with context in every case. The reason is that range reduction is
exact: $\theta$ is carried as a phase and reduced modulo $2\pi$ by truncation,
which is free and introduces no error, so the rotation applied at position
$m$ is no less accurate than the one applied at position $1$. Error is a
function of the angle's position within the reduced range, not of how large
$m$ has grown. The $F=9$ row for Qwen2 is non-monotonic, but that
configuration is already past the cliff of \S\ref{sec:precision} and we do not
read a trend into it.

This is a useful property for edge deployment: the precision chosen for a
short-context calibration remains valid as the context grows.
